%% kapa: exam=1 round=27 year=2016 period=2 form=A subject=appraisal_law subject_ko=감정평가관계법규 seq=24 number=24
%% source: Q-Net 감정평가사 기출문제 (cst00309) artlSeq=5169402 file=extracted/1008/5169402/제27회 감정평가사 1차 문제/제27회 감정평가사 2교시문제(A형).pdf
%% conversion: kordoc 4.24.51 -> tools/gen1.py (auto); review=proofread-sonnet; flags=none
\begin{question}{24}
건축법령상 건축허가에 관한 설명으로 옳은 것은?
\begin{choices}
  \item 위락시설에 해당하는 건축물의 건축을 허가하는 경우 건축물의 용도ㆍ규모가 주거환경 등 주변 환경을 고려할 때 부적합하다고 인정되면 건축위원회의 심의를 거쳐 건축허가를 하지 않을 수 있다.
  \item 연면적의 합계가 10만 제곱미터 이상인 공장을 광역시에 건축하려면 광역시장의 허가를 받아야 한다.
  \item 고속도로 통행료 징수시설을 대수선하려는 자는 시장ㆍ군수ㆍ구청장의 허가를 받아야 한다.
  \item 허가권자는 건축허가를 받은 자가 허가를 받은 날부터 6개월 이내에 공사에 착수하지 아니한 경우 허가를 취소하여야 한다.
  \item 건축위원회의 심의를 받은 자가 심의 결과를 통지 받은 날부터 1년 이내에 건축허가를 신청하지 아니하면 건축위원회 심의의 효력이 상실된다.
\end{choices}
\end{question}
\answer{1}
